% Generated draft: requires mathematical and source review.
\documentclass[11pt]{article}
\usepackage[margin=25mm]{geometry}
\usepackage{fontspec}
\setmainfont{Latin Modern Roman}
\usepackage{amsmath,amssymb,mathtools,unicode-math}
\setmathfont{Latin Modern Math}
\usepackage{enumitem,needspace,xurl,hyperref}
\setlength{\parindent}{0pt}
\setlength{\parskip}{5pt}
\setlength{\emergencystretch}{3em}
\newcommand{\N}{\mathbb{N}}
\newcommand{\Z}{\mathbb{Z}}
\newcommand{\Q}{\mathbb{Q}}
\newcommand{\R}{\mathbb{R}}
\newcommand{\C}{\mathbb{C}}
\newcommand{\F}{\mathbb{F}}
\newcommand{\A}{\mathbb{A}}
\newcommand{\PP}{\mathbb{P}}
\newcommand{\E}{\mathbb{E}}
\newcommand{\eps}{\varepsilon}
\newcommand{\dd}{\,\mathrm{d}}
\begin{document}
\section*{UnsolvedMath ID: 801157}
\textbf{Source problem code:} catalogue entry
\subsection{Definitions and mathematical statement}
Let \(m\in\N\) with \(m>1\). Write \(\mathcal G_m\) provisionally for the class called ``geometric Poisson structures of order \(m\)'' in the supplied source, and retain \(g^{ij}\) for the leading coefficient mentioned there. These are source-dependent designations: the supplied record does not define this class, the meaning of order, the underlying geometric setting, or the projective transformations under consideration.

The source poses the following questions:
\begin{enumerate}
\item For every integer \(m>1\), is every structure in \(\mathcal G_m\) projectively invariant?
\item Can the leading coefficient \(g^{ij}\) be interpreted as a projectively invariant geometric object?
\end{enumerate}

Here ``projectively invariant'' retains its unresolved source-specific meaning. The record supplies neither an expression containing \(g^{ij}\) nor its transformation law, the range of its indices, or the geometric type sought in the second question. Consequently, these questions preserve the supplied targets but require the missing source definitions before they constitute a fully specified mathematical formulation.
\subsection{Short English statement}
The problem asks whether every geometric Poisson structure of integer order greater than \(1\) is projectively invariant, and whether its leading coefficient \(g^{ij}\) has a projectively invariant interpretation. The supplied excerpt leaves the relevant structures and projective transformations undefined.
\subsection{Sources}
\begin{itemize}
\item Local UnsolvedMath registry, permanent ID \(801157\), source problem code \texttt{catalogue entry}, titled ``Nijenhuis Geometry and Integrable Systems --- Problem 10.'' The registry attributes the problem to Bolsinov, Konyaev, and Matveev, with proposed year \(2024\), and labels it open; this label is supplied metadata, not an independently verified mathematical assessment.
\item Bolsinov, Konyaev, and Matveev, \emph{Research Problems on Nijenhuis Geometry and Integrable Systems} (2024), Problem 10. Supplied source URL: \url{https://arxiv.org/abs/2410.04276}. The record reports access on \texttt{2026-10-16} and extraction by \texttt{arxiv\_latex}; these are supplied provenance claims, not independently verified access or source verification.
\end{itemize}
\end{document}
